% Copyright (c) 2026 Eugene Joseph M. Ragasa.
% SPDX-License-Identifier: Apache-2.0
% Author: Eugene Joseph M. Ragasa
% Formal proof completed: 2026-10-28
\documentclass[11pt]{article}
\usepackage[T1]{fontenc}\usepackage{amsmath,amssymb}\usepackage[margin=1in]{geometry}\usepackage[hidelinks]{hyperref}
\title{Proof Dossier: Rank-Two Swap Conjugation}\author{Eugene Joseph M. Ragasa}\date{Formal proof completed: 2026-10-28}
\begin{document}\raggedright\maketitle
\section{Identity and status}
ID \texttt{BG-FOURIER-THM-055}; status \texttt{proved}. Lean declaration:
\par\noindent{\tiny\path{ProjectKoios.BandGauge.Fourier.rankTwoSwapConjugation_firstCoordinateProjector}}.
\section{Statement and proof}
Let
\[
S=\begin{pmatrix}0&1\\1&0\end{pmatrix},\qquad P_0=\operatorname{diag}(1,0).
\]
Then $S^\dagger P_0S=\operatorname{diag}(0,1)$. Lean checks all four entries after expanding matrix multiplication.
\section{Boundary and audit}
This is exact rank-two algebra with no runtime or physical claim. Axioms are \texttt{propext},
\texttt{Classical.choice}, and \texttt{Quot.sound}; no admissions. Lean \texttt{4.35.0}, Mathlib
{\scriptsize\path{56a005855ea17e0b0958f73ec801ae140b0a1679}}; AI-assisted, kernel-checked, no independent expert review recorded.
\end{document}
